\documentclass[conference]{IEEEtran}
\IEEEoverridecommandlockouts

\usepackage{cite}
\usepackage{amsmath,amssymb,amsfonts}
\usepackage{graphicx}
\usepackage{textcomp}
\usepackage{xcolor}
\usepackage{booktabs}
\usepackage{multirow}
\usepackage{url}
\usepackage{hyperref}

\def\BibTeX{{\rm B\kern-.05em{\sc i\kern-.025em b}\kern-.08em
    T\kern-.1667em\lower.7ex\hbox{E}\kern-.125emX}}

\begin{document}

\title{Dissecting Positional Encoding in Encoder-Decoder MT:\\
A Cross-Lingual Study via Adaptive RoPE}

\author{
  \IEEEauthorblockN{Nilesh Mishra}
  \IEEEauthorblockA{
    Department of Computer Science and Engineering\\
    [Your Institution]\\
    thenileshmishra@gmail.com
  }
}

\maketitle

%% ─────────────────────────────────────────────────────────────────────────────
\begin{abstract}
Rotary Position Embedding (RoPE) has become the dominant positional encoding
in neural machine translation, yet its core assumption --- that all frequency
components contribute equally --- remains empirically untested in
encoder-decoder architectures. We introduce AdaptiveRoPE, a minimal extension
that augments RoPE with learnable per-head frequency gates and phase offsets.
Initialised to recover standard RoPE, these parameters reveal which
frequencies the model actually learns to use. We evaluate on three language
pairs (En--De with 3 seeds, Hi--En and Bn--En) against 5 baselines (RoPE,
Sinusoidal, ALiBi, and two ablations). Our findings are threefold:
(1)~trained gates consistently converge below one (mean~$\approx$0.90),
confirming a low-frequency preference;
(2)~the optimal adaptation is language-dependent: phase offsets are
critical for Hindi--En (+21\% BLEU over RoPE) while frequency gates matter
more for Bengali--En;
(3)~RoPE achieves the highest in-domain BLEU, but ALiBi degrades far less on
out-of-distribution lengths ($-$47\% vs.~$-$97\% at 201--400 words),
revealing a fundamental accuracy--extrapolation trade-off. These results
suggest that ``one-size-fits-all'' positional encodings are suboptimal and
motivate language-adaptive designs.
\end{abstract}

\begin{IEEEkeywords}
rotary position embedding, machine translation, positional encoding,
frequency analysis, cross-lingual study, encoder-decoder transformer
\end{IEEEkeywords}

%% ─────────────────────────────────────────────────────────────────────────────
\section{Introduction}
\label{sec:intro}

Position-aware attention is the backbone of sequence-to-sequence models.
Since the original sinusoidal encodings \cite{vaswani2017attention}, a
succession of relative and rotary approaches have improved how Transformers
encode position \cite{shaw2018self,raffel2020exploring,press2021train}.
Rotary Position Embedding (RoPE) \cite{su2021roformer} has emerged as the
de~facto standard: it rotates query and key vectors in 2-D frequency planes,
making the attention score depend only on relative position while preserving
the key--value cache and compilation compatibility. RoPE is now
used in LLaMA~\cite{touvron2023llama}, Mistral~\cite{jiang2023mistral}, and
state-of-the-art MT systems.

Despite its ubiquity, a critical assumption of RoPE has never been validated
in encoder-decoder MT: \emph{all frequency components receive equal weight}.
RoPE assigns each dimension pair $(2i, 2i{+}1)$ a fixed rotation frequency
$\omega_i = 10000^{-2i/d}$, with no mechanism to learn that some frequencies
are more useful than others. Recent work on decoder-only LLMs shows that
high-frequency RoPE dimensions can be pruned with minimal perplexity loss
\cite{men2024roundandround}, suggesting the uniform weighting is wasteful.
However, encoder-decoder MT differs fundamentally: positional information must
serve a bidirectional encoder, a causal decoder, \emph{and} cross-attention
between them. Whether the same frequency preference holds --- and whether it
varies across languages --- is unknown.

\textbf{Research questions.} This paper asks:
\begin{itemize}
  \item[\textbf{RQ1}] Do encoder-decoder MT models learn to suppress or amplify specific
        RoPE frequency components?
  \item[\textbf{RQ2}] Do the learned adaptations generalise across language pairs, or are
        positional encoding requirements language-specific?
  \item[\textbf{RQ3}] How do RoPE variants compare to alternatives (Sinusoidal, ALiBi) in
        terms of in-domain accuracy, length generalisation, and efficiency?
\end{itemize}

\textbf{Approach.} We design AdaptiveRoPE, a differentiable probe that adds
learnable per-head gates $\mathbf{g} \in \mathbb{R}^{H \times F}$ and phase
offsets $\boldsymbol{\phi} \in \mathbb{R}^{H \times F}$ to RoPE's rotation
angles (Section~\ref{sec:method}). Initialised at standard RoPE ($\mathbf{g}=\mathbf{1}$,
$\boldsymbol{\phi}=\mathbf{0}$), any divergence reflects a genuine learned
preference. Two ablations --- \textbf{GatesOnly} (phases fixed at zero) and
\textbf{PhasesOnly} (gates fixed at one) --- disentangle the contribution of
each component.

We train on WMT14 En--De (3 seeds), Samanantar Hi--En, and Samanantar Bn--En
with a 47M-parameter encoder-decoder Transformer. We evaluate 6 positional
encoding methods with BLEU, chrF, TER, and attention entropy; test
statistical significance; and measure length generalisation on 8 buckets up
to 400 source words.

\textbf{Key findings.}
\begin{itemize}
  \item Gates converge to $\approx$0.90 (Section~\ref{sec:gates}), confirming
        systematic low-frequency preference in encoder-decoder MT.
  \item \emph{Cross-lingual divergence:} PhasesOnly outperforms full
        AdaptiveRoPE on Hi--En (+21\% BLEU over RoPE), while AdaptiveRoPE
        itself wins on Bn--En (+0.8\% over RoPE). This reveals that the
        optimal PE strategy is language-dependent (Section~\ref{sec:crosslingual}).
  \item RoPE achieves the highest in-domain BLEU, but ALiBi generalises better
        to OOD lengths ($-$47\% vs.~$-$97\% degradation, Section~\ref{sec:length}).
  \item Attention entropy analysis shows ALiBi produces more uniform
        attention distributions (higher entropy), explaining its robustness
        to long inputs (Section~\ref{sec:attention}).
\end{itemize}

%% ─────────────────────────────────────────────────────────────────────────────
\section{Related Work}
\label{sec:related}

\textbf{Positional encodings.} Vaswani et al.~\cite{vaswani2017attention}
introduced sinusoidal absolute encodings. Shaw et al.~\cite{shaw2018self}
proposed relative position representations. ALiBi \cite{press2021train}
replaces explicit position encodings with a linear distance-based bias in
attention logits, achieving strong length extrapolation. RoPE
\cite{su2021roformer} uses rotation matrices to make dot products depend on
relative position; it dominates modern architectures due to cache efficiency
and compilation-friendly kernels.

\textbf{Analysing RoPE.} Men et al.~\cite{men2024roundandround} study
which RoPE frequencies are used in decoder-only Gemma-7B and find that the
highest frequencies can be truncated (p-RoPE). Our work differs in four ways:
(i)~we study \emph{encoder-decoder} MT, not decoder-only LLMs;
(ii)~we use a \emph{differentiable probe} rather than post-hoc ablation;
(iii)~we disentangle \emph{gates} from \emph{phases} via ablations;
(iv)~we evaluate across \emph{multiple language pairs} (En--De, Hi--En, Bn--En).

\textbf{Learnable RoPE variants.} YaRN \cite{peng2023yarn} and
LongRoPE \cite{ding2024longrope} learn frequency scalings for context-window
extension. These target length generalisation; we target understanding what
frequencies are needed at training-time lengths, with ablations that isolate
mechanisms.

\textbf{Cross-lingual PE studies.} Most PE research evaluates on English
or En--De \cite{su2021roformer,press2021train,men2024roundandround}. To our
knowledge, no prior work systematically compares PE requirements across
morphologically diverse language pairs with controlled architecture and training.

%% ─────────────────────────────────────────────────────────────────────────────
\section{Method}
\label{sec:method}

\subsection{Rotary Position Embedding (RoPE)}

For a query vector $\mathbf{q} \in \mathbb{R}^d$ at position $p$, RoPE
applies a block-diagonal rotation to each pair $(2i, 2i{+}1)$:
\begin{equation}
  \begin{bmatrix} q'_{2i} \\ q'_{2i+1} \end{bmatrix}
  =
  \begin{bmatrix} \cos\theta_p^{(i)} & -\sin\theta_p^{(i)} \\
                  \sin\theta_p^{(i)} &  \cos\theta_p^{(i)} \end{bmatrix}
  \begin{bmatrix} q_{2i} \\ q_{2i+1} \end{bmatrix},
\end{equation}
where $\theta_p^{(i)} = p \cdot \omega_i$ and $\omega_i = 10000^{-2i/d}$.
The same rotation is applied to keys, so the inner product depends only on
relative offset $p - p'$ \cite{su2021roformer}.

\subsection{Adaptive RoPE}

We augment the rotation angle with head-specific learnable gates
$\mathbf{g}$ and phase offsets $\boldsymbol{\phi}$:
\begin{equation}
  \theta_p^{(h,i)} = p \cdot \omega_i \cdot g^{(h,i)} + \phi^{(h,i)},
  \label{eq:adaptive}
\end{equation}
where $h \in \{1,\ldots,H\}$ indexes the attention head and $i \in
\{0,\ldots,F{-}1\}$ indexes the frequency ($F = d_{\text{head}}/2$).
Separate parameters are maintained for queries and keys:
$\mathbf{g}_q, \mathbf{g}_k, \boldsymbol{\phi}_q, \boldsymbol{\phi}_k \in
\mathbb{R}^{H \times F}$.

\textbf{Initialisation.} $\mathbf{g} = \mathbf{1}$, $\boldsymbol{\phi} =
\mathbf{0}$. This makes AdaptiveRoPE \emph{identical} to standard RoPE at
initialisation.

\textbf{Interpretation.} $g^{(h,i)} < 1$ suppresses frequency $i$;
$g^{(h,i)} > 1$ amplifies it. $\phi^{(h,i)}$ shifts the rotation angle
independently of frequency.

\textbf{Overhead.} For $H=6$, $F=32$, AdaptiveRoPE adds $4 \times 6 \times 32 =
768$ scalars per layer ($<0.002\%$ of 47M parameters).

\subsection{Ablations}

To isolate mechanisms we define two ablations:
\textbf{GatesOnly} has learnable $\mathbf{g}$ with $\boldsymbol{\phi}=\mathbf{0}$;
\textbf{PhasesOnly} has learnable $\boldsymbol{\phi}$ with $\mathbf{g}=\mathbf{1}$.
Comparing these to full AdaptiveRoPE and standard RoPE reveals whether frequency
scaling or phase shifting drives performance on each language pair.

%% ─────────────────────────────────────────────────────────────────────────────
\section{Experimental Setup}
\label{sec:setup}

\subsection{Datasets}

\textbf{WMT14 English--German} \cite{bojar2014wmt}: 4.5M training pairs,
\texttt{newstest2014} test (3,003 sentences). Tokenisation:
\texttt{Helsinki-NLP/opus-mt-en-de} SentencePiece (58K vocab).

\textbf{Samanantar Hi--En} and \textbf{Bn--En}
\cite{ramesh2022samanantar}: 1M and 450K training pairs respectively.
Test: 3,000 sentences sampled from official test splits.
Tokenisation: \texttt{ai4bharat/IndicBART} (64K vocab) \cite{dabre2021indicbart}.

\subsection{Model and Training}

Standard encoder-decoder Transformer: 6+6 layers, $d_{\text{model}}=384$,
$H=6$ heads, $d_{\text{ff}}=1536$, dropout 0.1, $\approx$47M parameters.
Embeddings tied with output projection.

Trained for 25K--30K steps, effective batch size 512, AdamW
\cite{loshchilov2018decoupled} with lr $10^{-3}$, cosine schedule, 5\%
warmup, gradient clipping 1.0, bf16 mixed precision, torch.compile,
label smoothing $\varepsilon=0.1$.

En--De uses seeds 42, 43, 44; Hi--En and Bn--En use seed 42.

\subsection{Evaluation Metrics}

\textbf{BLEU} \cite{papineni2002bleu}, \textbf{chrF}
\cite{popovic2015chrf}, and \textbf{TER} (Translation Edit Rate) via
SacreBLEU \cite{post2018sacrebleu}. Decoding: greedy (main results) and
beam search ($b=5$) for comparison.

\textbf{Attention analysis:} mean entropy and sink ratio (fraction of
attention mass on the first token) per layer, computed on 320 validation
samples.

\textbf{Length generalisation:} 8 source-length buckets (1--20, 21--40,
$\ldots$, 201--400 words). PE buffers extended to 420 tokens where needed.

\textbf{Statistical testing:} paired $t$-tests for multi-seed comparisons
(En--De); bootstrap resampling with 1,000 iterations for single-seed pairs.

%% ─────────────────────────────────────────────────────────────────────────────
\section{Results}
\label{sec:results}

\subsection{Translation Quality}
\label{sec:quality}

Table~\ref{tab:main} reports BLEU, chrF, TER, and inference speed for all
methods across three language pairs.

\textbf{En--De.} RoPE achieves the highest mean BLEU ($20.83 \pm 0.03$),
followed by PhasesOnly ($20.75 \pm 0.05$), AdaptiveRoPE ($20.69 \pm 0.05$),
and GatesOnly ($20.68 \pm 0.08$). All RoPE-based methods are tightly
clustered within 0.15 BLEU, but the differences are statistically
significant due to low inter-seed variance:
RoPE vs.~AdaptiveRoPE ($p = 0.001$), RoPE vs.~GatesOnly ($p = 0.002$),
and RoPE vs.~PhasesOnly ($p = 0.034$).
Sinusoidal ($20.18 \pm 0.15$) and ALiBi ($19.60$) lag by
$0.6$--$1.2$ BLEU, both highly significant ($p < 0.001$).

\textbf{Hi--En.} The ranking reverses dramatically: PhasesOnly achieves
$9.42$ BLEU, outperforming AdaptiveRoPE ($9.08$), GatesOnly ($8.37$), and
RoPE ($7.78$) --- a \textbf{21\% improvement} of PhasesOnly over RoPE.
This is our strongest cross-lingual finding: phase offsets are critical for
Hindi--English translation. Notably, full AdaptiveRoPE ($9.08$) underperforms
PhasesOnly ($9.42$), suggesting that the gates component interferes with
the beneficial phase adaptation on this language pair.

\textbf{Bn--En.} AdaptiveRoPE wins ($18.01$), followed by RoPE ($17.87$)
and GatesOnly ($17.79$). PhasesOnly drops to $17.06$, the lowest among
RoPE variants and $0.8$ BLEU below RoPE. This is the \emph{opposite} of
Hi--En: gates (frequency scaling) matter more than phases for Bengali.
The full AdaptiveRoPE outperforms both ablations, suggesting gates and
phases are complementary for this language pair.

\textbf{Speed.} ALiBi is fastest (55.6s) because it avoids RoPE's rotary
computation. Sinusoidal is similarly fast (54.7s). RoPE-based methods incur
a small overhead ($\approx$12--29\%) due to rotary kernel calls.
AdaptiveRoPE's overhead is negligible ($\approx$2s) relative to standard RoPE.

\textbf{Statistical significance.} Table~\ref{tab:sig} reports pairwise
$p$-values from Welch's $t$-tests on En--De BLEU scores. RoPE is
significantly better than all other methods except PhasesOnly (marginal,
$p=0.034$). Among RoPE variants, AdaptiveRoPE and GatesOnly are not
distinguishable ($p=0.835$), nor are AdaptiveRoPE and PhasesOnly
($p=0.162$). This confirms that while RoPE maintains a small but
consistent edge, the ablations cluster together.

\begin{table}[t]
\centering
\small
\caption{Pairwise statistical significance on En--De BLEU (Welch's $t$-test).
$^{**}p<0.01$, $^{*}p<0.05$.}
\label{tab:sig}
\begin{tabular}{lcccccc}
\toprule
& RoPE & ARoPE & Gates & Phases & Sinus & ALiBi \\
\midrule
RoPE    & --- & .001$^{**}$ & .002$^{**}$ & .034$^{*}$ & $<$.001$^{**}$ & N/A \\
ARoPE   & --- & --- & .835 & .162 & $<$.001$^{**}$ & N/A \\
Gates   & --- & --- & --- & .142 & $<$.001$^{**}$ & N/A \\
Phases  & --- & --- & --- & --- & $<$.001$^{**}$ & N/A \\
Sinus   & --- & --- & --- & --- & --- & N/A \\
\bottomrule
\end{tabular}
\end{table}

\begin{table*}[t]
\centering
\small
\caption{Translation quality and inference speed across positional encoding methods and language pairs.
En--De reports mean~$\pm$~std over 2--3 seeds; Hi--En and Bn--En use one seed. Bold = best per column.
$\uparrow$/$\downarrow$ indicate preferred direction.}
\label{tab:main}
\begin{tabular}{l|ccc|ccc|ccc|c}
\toprule
& \multicolumn{3}{c|}{\textbf{En--De}} & \multicolumn{3}{c|}{\textbf{Hi--En}} & \multicolumn{3}{c|}{\textbf{Bn--En}} & \\
\textbf{Method} & BLEU$\uparrow$ & chrF$\uparrow$ & TER$\downarrow$ & BLEU$\uparrow$ & chrF$\uparrow$ & TER$\downarrow$ & BLEU$\uparrow$ & chrF$\uparrow$ & TER$\downarrow$ & Decode~s$\downarrow$ \\
\midrule
RoPE & 20.83$\pm$0.03 & 52.51$\pm$0.19 & 68.20$\pm$0.06 & 7.78 & 38.29 & 158.74 & 17.87 & 42.70 & 79.56 & 62.0$\pm$0.3 \\
AdaptiveRoPE & 20.69$\pm$0.05 & 52.40$\pm$0.02 & 68.45$\pm$0.12 & 9.08 & 43.22 & 149.05 & \textbf{18.01} & 42.90 & 79.82 & 68.0$\pm$0.4 \\
GatesOnly & 20.68$\pm$0.08 & 52.45$\pm$0.08 & 68.48$\pm$0.19 & 8.37 & 41.43 & 169.07 & 17.79 & 42.97 & 79.57 & 67.6$\pm$0.3 \\
PhasesOnly & 20.75$\pm$0.05 & 52.46$\pm$0.11 & 68.47$\pm$0.03 & \textbf{9.42} & \textbf{43.26} & 175.90 & 17.06 & 42.16 & 83.16 & 71.6$\pm$6.0 \\
Sinusoidal & 20.18$\pm$0.15 & 52.01$\pm$0.06 & 69.06$\pm$0.08 & --- & --- & --- & 17.40 & 40.58 & \textbf{75.05} & 54.7$\pm$0.5 \\
ALiBi & 19.60 & 51.53 & 70.04 & --- & --- & --- & --- & --- & --- & \textbf{55.6} \\
\bottomrule
\end{tabular}
\end{table*}

\subsection{Cross-Lingual Ablations}
\label{sec:crosslingual}

The ablation results reveal language-dependent PE requirements
(Figure~\ref{fig:crosslingual}).

\textbf{Morphologically rich + free word order (Hindi).} PhasesOnly
$>$ AdaptiveRoPE $>$ GatesOnly $>$ RoPE. Phase offsets --- which shift
rotations without changing frequency scale --- are critical for capturing
Hindi's flexible syntax and rich morphology. RoPE's fixed phases are a
constraint the model strongly prefers to relax.

\textbf{Synthetic + relatively fixed order (Bengali).} AdaptiveRoPE
$>$ RoPE $\approx$ GatesOnly $>$ PhasesOnly. Frequency scaling (gates)
is more important than phase shifts. Bengali's more regular word order
and agglutinative morphology benefit from adaptive frequency weighting.

\textbf{Synthetic + rigid order (German).} All RoPE variants cluster
within 0.15 BLEU ($20.68$--$20.83$). Neither component dominates,
suggesting RoPE's defaults are near-optimal for Germanic syntax.

These patterns suggest that \emph{positional encoding requirements are not
universal}: morphological richness and word-order flexibility determine
whether phase or frequency adaptation is more valuable.

\begin{figure}[t]
  \centering
  \includegraphics[width=\linewidth]{../outputs/analysis/paper_artifacts_v2/crosslingual_comparison.png}
  \caption{BLEU comparison across language pairs. Error bars show std over
  seeds (En--De only).}
  \label{fig:crosslingual}
\end{figure}

\subsection{Per-Length Analysis}
\label{sec:perlength}

Table~\ref{tab:perlength} breaks BLEU down by source sentence length on
En--De. All methods perform best on mid-length sentences (11--50 words),
which dominate the training distribution. On very short sentences (1--10
words), RoPE leads by $0.3$--$0.7$ BLEU over ablations, suggesting that
its uniform frequency spectrum is well-suited for brief inputs. On long
sentences ($51+$ words, 21 samples), RoPE and AdaptiveRoPE tie
($\approx$22.6), with ALiBi close behind ($22.2$).

\begin{table}[b]
\centering
\scriptsize
\setlength{\tabcolsep}{3pt}
\caption{BLEU by source sentence length bucket on En--De (greedy decoding, mean over seeds).
Training distribution spans buckets 1--50 words.}
\label{tab:perlength}
\begin{tabular}{lrrrrr}
\toprule
\textbf{Method} & \textbf{1--10} & \textbf{11--20} & \textbf{21--30} & \textbf{31--50} & \textbf{51+} \\
\midrule
RoPE & 18.51$\pm$0.18 & 21.09$\pm$0.16 & 20.49$\pm$0.00 & 21.39$\pm$0.00 & 22.59$\pm$0.13 \\
ARoPE & 17.86$\pm$0.90 & 20.86$\pm$0.01 & 20.35$\pm$0.12 & 21.48$\pm$0.02 & 22.66$\pm$0.94 \\
Gates & 18.45$\pm$0.12 & 20.76$\pm$0.27 & 20.28$\pm$0.11 & 21.53$\pm$0.14 & 22.19$\pm$0.62 \\
Phases & 18.46$\pm$0.27 & 20.99$\pm$0.11 & 20.56$\pm$0.16 & 21.22$\pm$0.18 & 21.65$\pm$0.18 \\
Sinus & 17.94$\pm$0.36 & 20.12$\pm$0.31 & 19.88$\pm$0.02 & 21.07$\pm$0.35 & 21.76$\pm$0.42 \\
ALiBi & 17.77 & 19.63 & 19.28 & 20.23 & 22.23 \\
\bottomrule
\end{tabular}
\end{table}

\subsection{Length Generalisation}
\label{sec:length}

Table~\ref{tab:length} reports BLEU degradation on OOD source lengths
(En--De). ALiBi shows dramatically better extrapolation: only
$-$47\% degradation at 201--400 words, compared to $-$97\% for RoPE and
$-$98\% for AdaptiveRoPE variants. Sinusoidal also generalises better
($-$96\%) than RoPE-based methods.

Among RoPE variants, differences are small: GatesOnly ($-$31.8\% at
81--120 words) marginally outperforms RoPE ($-$36.7\%), but all converge
to near-zero BLEU at 201--400 words. The learned adaptations provide
insufficient extrapolation benefit to close the gap with ALiBi.

This reveals a fundamental \textbf{accuracy--extrapolation trade-off}:
RoPE achieves the highest in-domain BLEU but degrades most rapidly on
long sequences. ALiBi sacrifices $1.2$ BLEU in-domain for robust
extrapolation. For deployment with variable input lengths, the choice
between RoPE and ALiBi depends on the expected length distribution.

\begin{table}[b]
\centering
\scriptsize
\caption{BLEU degradation on OOD source lengths (En--De, greedy).
Values are \% change relative to the 1--20 word bucket.}
\label{tab:length}
\setlength{\tabcolsep}{3pt}
\begin{tabular}{lrrrrr}
\toprule
\textbf{Method} & \textbf{21--40} & \textbf{41--60} & \textbf{61--80} & \textbf{81--120} & \textbf{201--400} \\
\midrule
RoPE & +2.3 & +6.7 & $-$14.5 & $-$36.7 & $-$97.5 \\
ARoPE & +0.1 & +6.4 & $-$15.8 & $-$33.6 & $-$97.9 \\
Gates & +1.4 & +8.1 & $-$14.2 & $-$31.8 & $-$98.2 \\
Phases & $-$0.2 & +5.4 & $-$14.4 & $-$33.2 & $-$98.1 \\
Sinus & +5.1 & +10.4 & $-$11.9 & $-$30.9 & $-$95.8 \\
ALiBi & +1.3 & +6.1 & +0.3 & $-$17.0 & $-$94.8 \\
\bottomrule
\end{tabular}
\end{table}

\subsection{Attention Pattern Analysis}
\label{sec:attention}

Table~\ref{tab:attention} reports mean attention entropy and sink ratio.
ALiBi exhibits substantially higher entropy (2.233) than all other methods
(1.65--1.67), indicating more uniform attention distributions. Its higher
sink ratio (0.114 vs.~0.078 for RoPE) reflects stronger focus on initial
tokens --- a known ALiBi property that stabilises long-sequence attention
\cite{press2021train}.

Sinusoidal also shows elevated entropy and sink ratio, consistent with
its absolute encoding strategy. Among RoPE variants, entropy and sink
values are nearly identical, confirming that the learned gates and phases
do not fundamentally alter attention structure.

Figure~\ref{fig:entropy} shows per-layer entropy. All methods exhibit a
characteristic ``rise-and-fall'' pattern: entropy increases through the
mid-layers (layers 1--2) as contextual representations form, then
decreases in upper layers as attention sharpens. ALiBi maintains higher
entropy at every layer, explaining its robustness to length variation.

\begin{table*}[t]
\centering
\small
\caption{Attention entropy and sink ratio by PE method (En--De, mean across seeds).
Higher entropy = more uniform attention; higher sink = stronger focus on first token.}
\label{tab:attention}
\begin{tabular}{lcccccc}
\toprule
\textbf{Method} & RoPE & ARoPE & Gates & Phases & Sinus & ALiBi \\
\midrule
\textbf{Entropy}$\uparrow$ & 1.658$\pm$0.004 & 1.664$\pm$0.001 & 1.658$\pm$0.018 & 1.649$\pm$0.011 & 1.673 & 2.233 \\
\textbf{Sink}$\uparrow$ & 0.078$\pm$0.004 & 0.078$\pm$0.004 & 0.080$\pm$0.002 & 0.077$\pm$0.002 & 0.108 & 0.114 \\
\bottomrule
\end{tabular}
\end{table*}

\begin{figure}[t]
  \centering
  \includegraphics[width=0.9\linewidth]{../outputs/analysis/paper_artifacts_v2/entropy_per_layer.png}
  \caption{Per-layer attention entropy (En--De). ALiBi maintains higher
  entropy across all layers, producing more uniform attention.}
  \label{fig:entropy}
\end{figure}

%% ─────────────────────────────────────────────────────────────────────────────
\section{Analysis of Learned Parameters}
\label{sec:analysis}

\subsection{Learned Frequency Gates}
\label{sec:gates}

Table~\ref{tab:gates} reports gate statistics for En--De. All gates are
below one, confirming systematic suppression of the RoPE frequency
spectrum. Query gates (mean 0.915) are slightly higher than key gates
(0.903), indicating more aggressive frequency suppression in keys.

\begin{table}[t]
  \centering
  \small
  \caption{Learned gate and phase statistics (En--De, seed 43).
  Init: gates $=1.0$, phases $=0.0$.}
  \label{tab:gates}
  \begin{tabular}{lrrrr}
    \toprule
    \textbf{Parameter} & \textbf{Mean} & \textbf{Std} & \textbf{Min} & \textbf{Max} \\
    \midrule
    $\mathbf{g}_q$ (query gates)   & 0.9151 & 0.0249 & 0.8479 & 0.9744 \\
    $\mathbf{g}_k$ (key gates)     & 0.9030 & 0.0285 & 0.7784 & 0.9631 \\
    $\boldsymbol{\phi}_q$ (query phases) & 0.0045 & 0.0099 & $-$0.0216 & 0.0403 \\
    $\boldsymbol{\phi}_k$ (key phases)   & $-$0.0045 & 0.0099 & $-$0.0404 & 0.0217 \\
    \bottomrule
  \end{tabular}
\end{table}

\subsection{Frequency Spectrum Shift}

Figure~\ref{fig:spectrum} shows the effective frequency spectrum. The
learned model downscales effective frequencies across all dimensions, with
stronger suppression at higher-frequency indices. This means the model
prefers \emph{longer-range} positional signals, consistent with the need
to align source and target with potentially different word order.

\begin{figure}[htbp]
  \centering
  \includegraphics[width=0.9\linewidth]{../outputs/paper/fig_freq_spectrum.png}
  \caption{Effective frequency spectrum: standard RoPE (dashed) vs.\
  AdaptiveRoPE (solid). Shaded region = suppressed frequencies.}
  \label{fig:spectrum}
\end{figure}

\subsection{Per-Head and Per-Layer Specialisation}

Figure~\ref{fig:gates} visualises gate values per head and frequency.
Different heads learn distinct suppression profiles, confirming that
uniform frequency weighting is a constraint the model would relax. Head~2
shows the smallest deviation (least suppression), while head~0 shows the
largest, suggesting specialisation into long-range vs.\ short-range
positional attention.

Figure~\ref{fig:layers} shows per-layer encoder gates. Earlier layers
(0--2) suppress mid-to-high frequencies more aggressively, while later
layers (4--5) approach 1.0. This suggests the encoder progressively relies
more on positional information as representations become contextualised.

\begin{figure}[htbp]
  \centering
  \includegraphics[width=0.9\linewidth]{../outputs/analysis/adaptiverope_de/gates_q.png}
  \caption{Learned query gates $\mathbf{g}_q$ per head (rows) and
  frequency (columns). Red = suppression ($<1$); blue = amplification
  ($>1$).}
  \label{fig:gates}
\end{figure}

\begin{figure}[htbp]
  \centering
  \includegraphics[width=0.9\linewidth]{../outputs/paper/fig_gate_layers.png}
  \caption{Per-layer encoder gates, averaged over heads. Earlier layers
  suppress more aggressively.}
  \label{fig:layers}
\end{figure}

%% ─────────────────────────────────────────────────────────────────────────────
\section{Discussion}
\label{sec:discussion}

\textbf{Why are positional encoding requirements language-dependent?}
Our results suggest that typological properties --- morphological
complexity and word-order flexibility --- determine which PE components
matter. Hindi's rich fusional morphology and relatively free word order
benefit from \emph{phase offsets}, which provide flexible positional
shifts without altering the frequency spectrum. Bengali's more regular
agglutinative structure benefits from \emph{frequency scaling}, which
adjusts the positional ``resolution'' of different heads. German's rigid
SOV/V2 syntax is well-served by RoPE's defaults, leaving little room for
adaptation to improve.

\textbf{The accuracy--extrapolation trade-off.} RoPE's rotary structure
excels at in-domain lengths but fails catastrophically beyond the training
distribution ($-$97\% at 201--400 words). ALiBi's linear bias sacrifices
$1.2$ BLEU in-domain for dramatically better OOD performance ($-$47\%).
This trade-off is fundamental: RoPE's periodic encoding provides fine-grained
relative position discrimination, but the periodicity creates aliasing at
unseen lengths; ALiBi's monotonic bias generalises linearly but lacks the
fine-grained structure needed for peak in-domain performance.

\textbf{Implications for MT systems.} For high-resource language pairs
with controlled input lengths (e.g., document translation with length
filtering), RoPE remains the best choice. For low-resource or
length-variable settings (e.g., speech-to-text, live translation), ALiBi
or hybrid approaches may be preferable. Our cross-lingual findings suggest
that future work should explore \emph{language-conditional} positional
encodings that select gate/phase profiles based on typological features.

\textbf{Limitations.} Experiments are at a single scale (47M parameters).
While En--De uses 3 seeds, Hi--En and Bn--En use 1 seed due to
computational constraints. We evaluate greedy decoding; beam search may
reveal additional patterns. Whether findings generalise to larger models
and more language pairs remains open.

%% ─────────────────────────────────────────────────────────────────────────────
\section{Conclusion}
\label{sec:conclusion}

We presented AdaptiveRoPE, a minimal learnable extension of RoPE that reveals how
encoder-decoder MT models use the positional frequency spectrum. Through
experiments on three language pairs (En--De, Hi--En, Bn--En) against five
baselines, we established four key findings:
\begin{enumerate}
  \item Trained gates consistently converge below one
        (mean~$\approx 0.90$), confirming a systematic low-frequency
        preference in encoder-decoder MT.
  \item The optimal adaptation is language-dependent: phase offsets are
        critical for Hindi--En (+21\% BLEU over RoPE), while frequency
        gates matter more for Bengali--En.
  \item RoPE achieves the best in-domain BLEU, but ALiBi generalises
        better to OOD lengths ($-$47\% vs.~$-$97\% degradation).
  \item Attention entropy analysis reveals ALiBi's higher entropy
        explains its length robustness.
\end{enumerate}
These results challenge the ``one-size-fits-all'' assumption underlying
current positional encoding practice and motivate future work on
language-adaptive and task-adaptive PE designs.

%% ─────────────────────────────────────────────────────────────────────────────
\bibliographystyle{IEEEtran}
\bibliography{references}

\end{document}
